\chapter*{Abstract}
\addcontentsline{toc}{chapter}{Abstract}

Ho Chi Minh City operates a large public bus network whose Global Positioning System (GPS) telemetry and Automated Fare Collection (AFC) records produce a continuous stream of operational and demand data. The four existing public-facing tools available to residents (one route directory and three passenger routing applications) serve passenger navigation only; none of them provides analytical access to the network's performance, demand, or finance. The opening of the city's first metro line sharpens this gap by introducing an intermodal dependency between the bus network and the metro that the existing tools cannot monitor.

This report proposes an analytical and visualisation platform that integrates the two operational data streams behind a single queryable surface. The platform organises its analytical content into six themes (demand and crowding patterns, school-aware demand patterns, spatial and connectivity questions including bus-to-metro feeders, financial and operational efficiency, network-wide health and reliability, and per-route deep analysis). A four-layer warehouse stores raw views, conformed dimensions, fact views, and eight pre-computed aggregate tables, with the analytical surface reading exclusively from the aggregate layer to meet the interactive latency budget. The interaction model combines Multiple Coordinated Views (MCV) within each theme, an entity-detail interaction pattern on every clickable element, and the visual storytelling principles of Knaflic for chart presentation.

Phase One produces this report, four exploratory analysis notebooks (with field-level schema tables and five exploratory figures), a pixel-accurate user-interface mockup, and a partial implementation of the data-preprocessing pipeline at Layer 1 of the proposed architecture. The exploratory analysis finds that student fares account for approximately 34 percent of weekday-morning boardings, cash dominates payments at approximately 96 percent, route demand is heavy-tailed with the top 55 routes carrying 80 percent of total volume, and the Tết period produces a visible nine-day ridership trough.

The Phase Two implementation plan completes the warehouse and the analytical surface across nine sequential phases of approximately twenty-one working days, followed by a two-wave user evaluation across the four user roles. Three extensions are recorded as candidates for follow-on work: real-time arrival-time prediction, per-route anomaly detection, and bus telemetry as a city-wide traffic-sensing substrate.

\vspace{1em}
\noindent\textbf{Keywords:} public transit, bus operational analytics, visual analytics, multiple coordinated views, data warehouse.
